\documentclass[11pt]{article}
\usepackage[margin=1in]{geometry}
\usepackage[hidelinks]{hyperref}
\input{numbers.tex}
\begin{document}

\noindent \today \\[1em]
The Editor-in-Chief \\
\textit{Electric Power Systems Research} \\
Elsevier \\[1em]
Dear Editor,

\medskip
Please consider our manuscript \textbf{``Schedule-Aware Real-time Depot Charging for Electric Fleets''} by
Mohammad Kamar Uddin, Ali Moradi Amani and Hui Song for publication as a research paper in
\textit{Electric Power Systems Research}.

Commercial fleet depots charge more vehicles than they have chargers, under a transformer limit, while each
vehicle must leave with enough energy for a rostered trip that can change during the day. The manuscript
proposes AQPS, a three-layer controller that (i) derives each vehicle's mission energy from the roster and
admits the largest priority set for which a convex receding-horizon quadratic program can meet every
mission as a hard constraint, (ii) allocates chargers with just-in-time preemption and a safe-release rule
that frees a charger from a priority vehicle only after its mission energy is delivered, and (iii) solves
the quadratic program only on events, feeding predicted shortfalls back to admission, with an operator
impact preview for disruptions. We prove that the problem is convex with a unique optimal aggregate profile
and give sufficient conditions under which every admitted priority vehicle departs with its mission energy.

The evaluation uses real AEMO rooftop-PV forecasts and actuals, seven methods including a
perfect-information mixed-integer benchmark, scripted and random disruptions, overload and sensitivity
studies. In the peak-stress scenario AQPS met all missions, against \RMissionLLFSThreeFortyFive\,\% and
\RMissionLLFSThreeNinety\,\% for least-laxity-first, at an equal or lower cost per kWh, and event triggering
cut optimiser solves by \RTotRedMin--\RTotRedMax\,\%. We also report where simple heuristics are as good and
the reduction in non-priority energy that the mission guarantee costs. Every number in the paper is
generated by scripts from the committed simulation results.

% Optional (authors' decision): disclose the earlier submission.
An earlier version of this work was submitted to the \textit{Journal of Energy Storage} (EST-D-26-06112) and
was not accepted. The present manuscript has been substantially revised in response to those reviews: the
priority layer, the safe-release rule, the event-triggered solver, the PV data (now real AEMO data), the
baselines and all results are new.

\textbf{Article length.} The manuscript is a full-length article of about 6{,}000 words (abstract, text,
tables and figure captions; references excluded), or about 6{,}600 words if the text inside figures is
counted, within the 7{,}000-word limit.

The manuscript is original, has not been published elsewhere and is not under consideration by another
journal. All authors have approved it and declare no competing interests.

\medskip
\noindent Yours sincerely, \\[1em]
\textbf{Mohammad Kamar Uddin} (corresponding author) \\
School of Engineering, RMIT University, Melbourne, VIC 3000, Australia \\
\texttt{s3955218@student.rmit.edu.au} \\[0.5em]
on behalf of Ali Moradi Amani (RMIT University) and Hui Song (University of Newcastle)

\end{document}
